\chapter{Threads, Resource Allocation, and Deadlocks}\label{ch:deadlock}
\chapterguide{Chapter 5: shared state, race conditions, and synchronization.}{distinguish processes from threads, use a mutex to protect a critical section, identify the four deadlock conditions, and reason about safe states in Lab 6.}

\section{POSIX threads}
The POSIX threads API creates multiple execution streams inside one process. \texttt{pthread\_create()} starts a thread, \texttt{pthread\_self()} identifies the calling thread, \texttt{pthread\_equal()} compares thread identifiers, and \texttt{pthread\_join()} waits for completion.

\begin{lstlisting}[style=utpc]
pthread_t t;
int err = pthread_create(&t, NULL, worker, NULL);
if (err != 0) {
    /* handle error */
}
pthread_join(t, NULL);
\end{lstlisting}

\begin{pitfall}
Thread identifiers are opaque values. Portable code compares them with \texttt{pthread\_equal()} rather than assuming they are ordinary integers with a specific representation.
\end{pitfall}

\section{Mutexes}
A \term{mutex} protects a critical section so only one thread owns the lock at a time. The normal lifecycle is initialize, lock, access shared state, unlock, then destroy after all users are finished.

\begin{utpfig}[What threads share]{What the threads of one process share. Code, globals, and the heap exist once and are seen by every thread; each thread has its own stack and registers, which need no synchronization.}{fig:thread-memory}
\begin{tikzpicture}[
  seg/.style={draw=UTPBlue,fill=UTPPaleBlue,minimum width=3.0cm,minimum height=0.8cm,
              font=\scriptsize\color{UTPNavy}},
  shared/.style={seg,fill=UTPPaleOrange,draw=UTPOrange},
  stk/.style={draw=UTPTeal,fill=UTPPaleTeal,minimum width=1.7cm,minimum height=0.8cm,
              font=\scriptsize\color{UTPNavy}},
  lab/.style={font=\scriptsize\color{UTPMuted},align=center},
  warn/.style={font=\scriptsize\color{UTPOrange},align=center}]

\draw[draw=UTPNavy,rounded corners=2mm,line width=0.9pt] (-2.1,-2.75) rectangle (8.7,1.75);
\node[font=\small\bfseries\color{UTPNavy},anchor=north west] at (-2.0,1.65) {one process};

\node[seg]    at (-0.4,0.85) {code (text)};
\node[shared] at (-0.4,-0.35) {globals};
\node[shared] at (-0.4,-1.55) {heap};
\node[lab,text width=2.6cm] at (-0.4,-2.30) {one copy, seen by every thread};

\foreach \i/\x in {1/2.9, 2/5.2, 3/7.5} {
  \node[stk] at (\x,0.85) {stack \i};
  \node[stk] at (\x,-0.35) {registers \i};
  \node[lab] at (\x,-1.08) {thread \i};
}
\node[lab,text width=5.6cm] at (5.2,-2.05) {private to each thread --- no synchronisation needed};

\node[warn,text width=11cm] at (3.3,-3.55)
  {The orange region is the whole reason a mutex exists. Two threads writing a
   global without one is a data race, and the outcome depends on scheduling ---
   which is why such a bug reproduces only sometimes.};

\end{tikzpicture}
\end{utpfig}

\section{What deadlock means}
A \term{deadlock} is a state in which a set of processes or threads are permanently blocked because each waits for an event that only another member of the set can cause.

Four Coffman conditions must hold simultaneously for resource deadlock:
\begin{enumerate}
  \item \textbf{Mutual exclusion}: at least one resource cannot be shared concurrently.
  \item \textbf{Hold and wait}: a process holds resources while requesting more.
  \item \textbf{No preemption}: resources cannot simply be taken away.
  \item \textbf{Circular wait}: a cycle of processes each waits for a resource held by the next.
\end{enumerate}
Breaking any one condition prevents this class of deadlock.

\begin{utpfig}[A circular wait]{A circular wait. P1 holds R1 and requests R2, while P2 holds R2 and requests R1, so neither can proceed.}{fig:circular-wait}
\begin{tikzpicture}[
  proc/.style={circle,draw=UTPBlue,fill=UTPPaleBlue,minimum size=1.05cm,
               font=\small\bfseries\color{UTPNavy}},
  res/.style={draw=UTPNavy,fill=white,minimum width=1.05cm,minimum height=0.9cm,
              font=\small\bfseries\color{UTPNavy}},
  cyc/.style={-{Stealth[length=2.3mm]},line width=1pt,draw=UTPRed},
  lab/.style={font=\scriptsize\color{UTPMuted}},
  note/.style={font=\scriptsize\color{UTPInk},align=center}]

\node[proc] (p1) at (0,2.5)   {P1};
\node[res]  (r2) at (5.0,2.5) {R2};
\node[proc] (p2) at (5.0,0)   {P2};
\node[res]  (r1) at (0,0)     {R1};

\draw[cyc] (p1) -- node[lab,above]{requests} (r2);
\draw[cyc] (r2) -- node[lab,right]{assigned to} (p2);
\draw[cyc] (p2) -- node[lab,below]{requests} (r1);
\draw[cyc] (r1) -- node[lab,left]{assigned to} (p1);

\node[note,text width=4.2cm] at (2.5,1.25)
  {the cycle \emph{is} the circular wait --- neither process can move, and
   nothing outside the cycle will move them};

\node[note,text width=11cm,text=UTPMuted] at (2.5,-1.15)
  {Each resource has one instance, so this cycle is a deadlock. Remove any one
   arrow --- by preemption, by releasing before requesting, or by ordering the
   resources --- and the cycle opens.};

\end{tikzpicture}
\end{utpfig}

\section{Prevention, avoidance, and detection}
\begin{utptable}
\renewcommand{\arraystretch}{1.18}
\begin{tabularx}{\linewidth}{@{}>{\bfseries\leavevmode\color{UTPNavy}}m{0.22\linewidth}
  >{\raggedright\arraybackslash}X
  >{\raggedright\arraybackslash}X@{}}
\toprule
\rowcolor{UTPPaleBlue!92}
\utpth{Strategy} & \utpth{Core idea} & \utpth{Trade-off} \\
\midrule
Prevention & Design the protocol so a Coffman condition cannot hold & Can reduce utilization or flexibility \\
Avoidance & Grant a request only if the resulting state remains safe & Requires future maximum-demand information \\
Detection/recovery & Allow deadlock, detect it later, then recover & Detection and recovery can be expensive \\
\bottomrule
\end{tabularx}
\end{utptable}

\section{Banker's algorithm}
Banker's algorithm avoids unsafe allocations. First calculate
\texttt{Need = Max - Allocation}. Start with \texttt{Work = Available}, then
find a process whose needs fit within Work. Assume it finishes and returns its
allocation to Work. If every process can finish in some order, the state is
\term{safe}.

\begin{workedexample}
Suppose Available is $(1,1)$, P1 has Need $(1,0)$ and Allocation $(1,1)$, while P2 has Need $(2,1)$ and Allocation $(1,0)$. P1 can finish first because $(1,0)\le(1,1)$. Work becomes $(2,2)$, allowing P2 to finish. Therefore a safe sequence is P1, P2.
\end{workedexample}

\begin{utpfig}[The Banker's safety check]{The Banker's safety check for the worked example. P1's need fits within Work, so it can finish and release its allocation; the larger Work then lets P2 finish, so P1, P2 is a safe sequence.}{fig:bankers}
\begin{tikzpicture}[
  wk/.style={draw=UTPTeal,fill=UTPPaleTeal,rounded corners=1.5mm,minimum width=1.5cm,
             minimum height=0.75cm,font=\small\bfseries\color{UTPNavy}},
  test/.style={draw=UTPBlue,fill=UTPPaleBlue,rounded corners=1.5mm,text width=3.0cm,
               minimum height=1.15cm,align=center,font=\scriptsize\color{UTPNavy}},
  go/.style={-{Stealth[length=2.2mm]},thick,draw=UTPTeal},
  wcap/.style={font=\scriptsize\bfseries\color{UTPMuted}},
  rel/.style={font=\scriptsize\color{UTPGreen},align=center}]

\node[wcap] at (0,1.15) {Work = Available};
\node[wk] (w0) at (0,0.35) {(1,\,1)};

\node[test] (t1) at (3.5,0.35) {P1 needs (1,\,0)\\$(1,0)\le(1,1)$ --- P1 can finish};
\node[wk]   (w1) at (7.0,0.35) {(2,\,2)};
\node[rel]  at (7.0,-0.62) {P1 releases (1,\,1)};

\node[test] (t2) at (10.5,0.35) {P2 needs (2,\,1)\\$(2,1)\le(2,2)$ --- P2 can finish};
\node[wk]   (w2) at (14.0,0.35) {(3,\,2)};

\draw[go] (w0) -- (t1);
\draw[go] (t1) -- (w1);
\draw[go] (w1) -- (t2);
\draw[go] (t2) -- (w2);

\node[font=\scriptsize\color{UTPInk},align=center,text width=13cm] at (7.0,-1.55)
  {Every process fits into the order, so the state is \textbf{safe} and the
   safe sequence is P1, P2. Had P1 failed its test as well, no order would
   exist and the request would be refused --- not because deadlock had
   happened, but because it could no longer be ruled out.};

\end{tikzpicture}
\end{utpfig}

\begin{pitfall}
An \term{unsafe state} is not identical to an already existing deadlock. Unsafe means the OS can no longer guarantee that every future maximum claim can be satisfied in some order. Deadlock is an actual circular waiting condition.
\end{pitfall}

\section{Detection versus the supplied exercise}
A textbook deadlock-detection algorithm normally uses current outstanding requests rather than each process's declared future maximum. The supplied lab calculates \texttt{Need = Max - Allocation}; this is close to a safety/avoidance test. Treat the code as a useful matrix exercise, but keep the conceptual distinction clear.

\labsection{6}{Threads, mutexes, and resource safety}{sec:lab06}
\begin{labproject}{See a race condition, protect it, then trace a safe-state test}

\textbf{Coverage.} Chapter~\ref{ch:deadlock}: threads, critical sections, mutexes, safe resource states.\\
\textbf{Goal.} Run the same increment loop with and without a lock, then inspect a Banker trace.\\
\textbf{Explain.} Explain why a race is timing-dependent and why ``unsafe'' is not identical to ``deadlocked now''.

\section*{Source code}
\labsourcecard{\texttt{threads\_demo.c}}{Two threads increment one shared counter.}{The read-increment-write sequence is protected only when the mutex is enabled.}
\lstinputlisting[style=utpcnumbered]{all_experiments/Lab06/threads_demo.c}

\labsourcecard{\texttt{banker.c}}{Compute \texttt{Need = Maximum - Allocation} and search for a safe sequence.}{A process may finish only when every resource in its Need vector is no larger than current Work.}
\lstinputlisting[style=utpcnumbered]{all_experiments/Lab06/banker.c}

\section*{Task 1 --- Repeat the race, do not trust one run}
\begin{lstlisting}[style=utpbash]
cd all_experiments/Lab06
gcc -Wall -Wextra -std=c11 -pthread threads_demo.c -o threads_demo
for i in 1 2 3 4 5; do ./threads_demo --no-lock; done
\end{lstlisting}
Record all five Actual values. A race condition means the result is not protected by a correctness guarantee; it does \emph{not} mean every run must fail differently.

\section*{Task 2 --- Enable the mutex}
\begin{lstlisting}[style=utpbash]
for i in 1 2 3; do ./threads_demo; done
\end{lstlisting}
The locked runs should reach the Expected count. On paper, mark the critical section: read counter $\rightarrow$ increment $\rightarrow$ write counter.

\section*{Task 3 --- Safe and unsafe resource states}
\begin{lstlisting}[style=utpbash]
gcc -Wall -Wextra -std=c11 banker.c -o banker
./banker
./banker --unsafe
\end{lstlisting}
For the safe run, copy the sequence and verify the first selected process by comparing its Need with the initial Work vector. For the unsafe run, identify the first point at which no unfinished process can satisfy \texttt{Need <= Work}.

\begin{utptable}
\renewcommand{\arraystretch}{1.18}
\begin{tabularx}{\linewidth}{@{}>{\bfseries\leavevmode\color{UTPNavy}}m{0.38\linewidth}
  >{\raggedright\arraybackslash}X@{}}
\toprule
\rowcolor{UTPPaleBlue!92}
\utpth{Observation} & \utpth{Interpretation} \\
\midrule
Unlocked counter loses updates & timing allowed two threads to overlap a non-atomic update \\
Locked counter reaches expected value & one thread at a time enters the protected update \\
Banker finds full sequence & there exists an order in which all declared needs can complete \\
Banker cannot find full sequence & the state is unsafe for the declared maximum claims \\
\bottomrule
\end{tabularx}
\end{utptable}

\begin{quickcheck}
\begin{enumerate}
  \item What memory do threads in one process share?
  \item What does a mutex guarantee when used correctly around a critical section?
  \item What is the difference between a deadlocked state and an unsafe state?
\end{enumerate}
\end{quickcheck}

\begin{verification}
\begin{itemize}
  \item Five unlocked runs are recorded, not one.
  \item All three locked runs reach the Expected count.
  \item The critical section is marked as read, increment, write.
  \item The Banker's check identifies one safe and one unsafe state, and an unsafe state is distinguished from an actual deadlock.
\end{itemize}
\end{verification}

\begin{labdeliverable}
Submit the five unlocked values, three locked values, one marked critical-section sketch, and safe/unsafe Banker evidence. Include a short distinction between \textbf{unsafe state} and \textbf{actual deadlock}.
\end{labdeliverable}
\end{labproject}

